\documentclass[a4paper]{article}

\usepackage[fontset=fandol]{ctex}
\usepackage{amsmath, amssymb}
\usepackage{graphicx}
\usepackage[margin=2.5cm]{geometry}
\usepackage[most]{tcolorbox}
\usepackage{listings}
\usepackage{hyperref}
\usepackage{booktabs}
\usepackage{float}
\usepackage{tikz}

\setlength{\parindent}{2em}
\setlength{\parskip}{0.35em}
\sloppy

\newtcolorbox{knowledgebox}[1]{
    enhanced, colback=blue!5!white, colframe=blue!75!black, colbacktitle=blue!75!black,
    coltitle=white, fonttitle=\bfseries, title={#1}, attach boxed title to top left={yshift=-2mm, xshift=2mm},
    boxrule=1pt, sharp corners
}

\newtcolorbox{importantbox}[1]{
    enhanced, colback=yellow!10!white, colframe=yellow!80!black, colbacktitle=yellow!80!black,
    coltitle=black, fonttitle=\bfseries, title={#1}, sharp corners
}

\newtcolorbox{warningbox}[1]{
    enhanced, colback=red!5!white, colframe=red!75!black, colbacktitle=red!75!black,
    coltitle=white, fonttitle=\bfseries, title={#1}, sharp corners
}

\lstset{
    basicstyle=\ttfamily\small,
    keywordstyle=\color{blue},
    stringstyle=\color{red!60!black},
    commentstyle=\color{green!60!black},
    breaklines=true,
    frame=single,
    numbers=left,
    numberstyle=\tiny\color{gray},
    captionpos=b,
    extendedchars=false
}

\newcommand{\notetitle}{Transformer（下）：Decoder、Cross Attention 与训练技巧}
\newcommand{\notesubtitle}{李宏毅《机器学习》2025 第六节：从自回归解码到 exposure bias}
\newcommand{\noteauthors}{基于李宏毅老师课程整理}
\newcommand{\notedate}{\today}
\newcommand{\videochannel}{李宏毅机器学习深度学习课程（Bilibili）}
\newcommand{\videoduration}{约 61 分钟}
\newcommand{\videourl}{https://www.bilibili.com/video/BV1TAtwzTE1S?p=16}

\begin{document}
\pagestyle{empty}

\begin{center}
    \vspace*{1.5cm}
    {\Huge\bfseries \notetitle\par}
    \vspace{0.4cm}
    {\LARGE \notesubtitle\par}
    \vspace{0.8cm}
    \includegraphics[width=0.62\textwidth]{video.jpg}\par
    \vspace{0.8cm}
    {\large \noteauthors\par}
    {\large \notedate\par}
    \vspace{0.8cm}
    \begin{tcolorbox}[width=0.82\textwidth,center,sharp corners,
                      colback=black!3!white,colframe=black!50!white]
        \centering
        \textbf{视频信息}\\[3pt]
        频道：\videochannel \\
        时长：\videoduration \\
        链接：\href{\videourl}{\videourl}
    \end{tcolorbox}
\end{center}

\newpage
\tableofcontents
\newpage
\pagestyle{plain}

% =====================================================================
\section{Decoder 的任务：把表示变成序列}
% =====================================================================

上一讲已经把 Transformer encoder 拆开：encoder 接收输入序列，输出一排上下文化向量。本讲继续讲 decoder。Seq2seq 的完整流程可以写成
\[
    X \xrightarrow{\mathrm{Encoder}} H
    \xrightarrow{\mathrm{Decoder}} Y,
\]
其中 $X$ 是输入序列，$H$ 是 encoder 输出，$Y=(y_1,\ldots,y_N)$ 是 decoder 产生的输出序列。

Decoder 要解决的问题比“分类一次”复杂得多。它不仅要决定每一步输出哪个 token，还要决定什么时候停止。以语音辨识为例，输入是一段声音，输出可能是中文字序列“机器学习”。模型不能一开始就知道输出一定是 4 个字；它必须边产生、边根据已经产生的前缀继续预测。

\subsection{输出单位：字、subword 或 token}

Decoder 的输出单位可以有很多选择。英文常用 subword、wordpiece、BPE token 或 byte-level token；中文课程里为了说明方便，常把一个中文字当作一个 token。实际系统会把所有可能输出组成 vocabulary，大小记作 $V$。每一步 decoder 输出一个长度为 $V$ 的分布：
\[
    p_t = P(y_t \mid y_{<t}, H).
\]
这里 $p_t$ 的每个维度对应一个候选 token。模型可以取最大概率 token，也可以从分布中 sampling；训练时则用正确答案监督这个分布。

\begin{importantbox}{Decoder 的三个基本问题}
    Decoder 必须回答三件事：第一，每一步该看哪些过去 token；第二，每一步如何从 encoder 输出抽取相关信息；第三，输出序列何时结束。这三件事分别对应 masked self-attention、cross attention 和 END token。
\end{importantbox}

\subsection{本章小结}

Decoder 把 encoder 表示转成目标序列。它的输出不是一次性固定维度向量，而是一串 token 分布；每一步都依赖 encoder 信息和已经生成的前缀。

% =====================================================================
\section{Autoregressive Decoder}
% =====================================================================

最常见的 decoder 是 autoregressive decoder，简称 AT decoder。它一次产生一个 token：先读入特殊的 BEGIN token，产生第一个字；再把第一个字当作下一步输入，产生第二个字；如此循环，直到输出 END token。

用概率写就是
\[
    P(Y\mid X)
    =
    \prod_{t=1}^{N}
    P(y_t\mid y_1,\ldots,y_{t-1},H).
\]
其中 $H=\mathrm{Encoder}(X)$。这个分解强调了“前缀条件”：生成第 $t$ 个 token 时，模型只应该依赖已经生成的 $y_1,\ldots,y_{t-1}$，不能看未来答案。

\begin{figure}[H]
    \centering
    \includegraphics[width=0.88\textwidth]{figures/fig01_autoregressive_decoder.jpg}
    \caption{Autoregressive decoder 会把上一轮输出接回下一轮输入；图中红色虚线表示模型生成“机、器、学”后，再继续产生“习”。\protect\footnotemark}
    \label{fig:at-decoder}
\end{figure}
\footnotetext{视频画面时间区间：约 08:20--08:50；字幕对应“decoder 看到自己产生出来的错误输入”“error propagation”“一步错步步错”。}

\subsection{Error Propagation}

AT decoder 的优势是自然地处理可变长度和前缀依赖；问题是错误会传播。若模型第一步本来应该输出“机”，却输出成“气”，第二步看到的前缀就已经错了。之后模型是在错误上下文上继续预测，可能越走越偏。

这也是后面 exposure bias 和 scheduled sampling 会出现的原因。训练时模型常常看到正确前缀；测试时却只能看到自己的预测前缀。一旦预测错，模型未必知道如何从错误状态恢复。

\begin{warningbox}{不要把 AT 解码误解成一次输出整句}
    推论时，AT decoder 不是同时知道所有输出 token。它先产生第一个，再把第一个放回输入产生第二个。训练图里常画成并行输出，是因为 teacher forcing 可以把正确前缀一次喂进去，但这不是普通推论时的工作方式。
\end{warningbox}

\subsection{本章小结}

Autoregressive decoder 用链式法则把整句概率拆成逐步生成。它适合可变长度输出，也能利用已经生成的上下文，但会带来错误传播和训练/推论不一致的问题。

% =====================================================================
\section{Masked Self-attention 与 END Token}
% =====================================================================

Transformer decoder 里的第一块 attention 不是普通 self-attention，而是 masked self-attention。普通 self-attention 允许每个位置看所有位置；decoder 生成第 $t$ 个 token 时，未来 token 还不存在，因此位置 $t$ 不能看 $t+1,t+2,\ldots$。

在实现上，可以给 attention score 加一个 causal mask：
\[
    e_{ij}=\frac{q_i^\top k_j}{\sqrt{d_k}},
    \qquad
    \alpha_{ij}
    =
    \mathrm{softmax}_j(e_{ij}+M_{ij}),
\]
其中
\[
    M_{ij}=
    \begin{cases}
        0, & j\le i,\\
        -\infty, & j>i.
    \end{cases}
\]
于是第 $i$ 个位置只能 attend 到自己和左边的位置。

\begin{figure}[H]
    \centering
    \includegraphics[width=0.82\textwidth]{figures/fig02_masked_self_attention.jpg}
    \caption{Masked self-attention 中，计算 $b^2$ 时只能使用 $a^1,a^2$，不能使用还没生成的 $a^3,a^4$。\protect\footnotemark}
    \label{fig:masked}
\end{figure}
\footnotetext{视频画面时间区间：约 12:50--13:20；字幕对应“计算 b2 的时候没有 A3 跟 A4”“decoder 输出是一个一个产生的，所以只能考虑左边”。}

\subsection{为什么需要 END token}

Decoder 还要决定什么时候停止。做法是在 vocabulary 里加入一个特殊 token，例如 END。模型每一步都输出一个分布，END 和普通 token 一样占一个维度。一旦模型输出 END，生成过程结束。

\begin{figure}[H]
    \centering
    \includegraphics[width=0.82\textwidth]{figures/fig03_stop_token.jpg}
    \caption{除了普通中文字，decoder 的输出分布里还要有 END 这样的特殊符号；模型输出 END 时，序列生成停止。\protect\footnotemark}
    \label{fig:end-token}
\end{figure}
\footnotetext{视频画面时间区间：约 15:25--16:05；字幕对应“准备一个特别的符号”“用 END 来表示”“开始跟段可以用同一个符号表示”。}

BEGIN 和 END 有时可以使用同一个符号，因为 BEGIN 只出现在 decoder 输入端，END 只出现在输出端；也可以使用两个不同符号。关键不是符号名字，而是训练资料必须告诉模型：什么时候应该开始，什么时候应该停止。

\subsection{本章小结}

Masked self-attention 保证 decoder 不能偷看未来；END token 让 decoder 能自己决定输出长度。两者共同把 Transformer decoder 变成真正可生成序列的模块。

% =====================================================================
\section{AT 与 NAT：逐步生成还是平行生成}
% =====================================================================

Autoregressive decoder 一次产生一个 token，优点是建模自然，缺点是推论慢，因为第 $t$ 步必须等第 $t-1$ 步完成。Non-autoregressive decoder，简称 NAT decoder，尝试一次并行地产生多个 token。它的典型输入可能是多个 BEGIN token，然后同时输出 $w_1,w_2,\ldots$。

\begin{figure}[H]
    \centering
    \includegraphics[width=0.86\textwidth]{figures/fig04_at_vs_nat.jpg}
    \caption{AT decoder 逐步生成并自然产生 END；NAT decoder 平行输出，需要额外处理输出长度，并常遇到 multi-modality 问题。\protect\footnotemark}
    \label{fig:at-nat}
\end{figure}
\footnotetext{视频画面时间区间：约 22:50--23:35；字幕对应“NAT performance 往往不如 AT”“需要很多 trick 才能逼近 AT”“multi-modality 的问题”。}

\subsection{NAT 的输出长度问题}

NAT 不能像 AT 那样自然地生成到 END 为止，因为它试图同时输出。常见解决方式有两种：
\begin{itemize}
    \item 另外训练一个 length predictor，先预测输出长度，再让 NAT 产生对应数量 token；
    \item 输出一个很长的序列，遇到 END 后忽略后面 token。
\end{itemize}

NAT 的优势是平行化，推论速度可能更快，输出长度也更容易控制。但它通常比 AT 难训练，效果往往较差。

\subsection{Multi-modality 直觉}

NAT 的困难之一是 multi-modality。假设一句话有多个合理翻译，AT decoder 可以在第一步选择某个方向，然后后续 token 跟着这个方向保持一致。NAT 同时预测所有位置，每个位置可能各自选择不同模式，导致组合起来不一致。例如有的位置像第一种翻译，有的位置像第二种翻译，最终句子不自然。

\begin{knowledgebox}{为什么 AT 容易保持一致}
    AT 每一步都把已经生成的前缀作为条件，因此早期选择会约束后续选择。NAT 平行生成时缺少这种逐步承诺机制，所以需要额外技巧让不同位置之间协调。
\end{knowledgebox}

\subsection{本章小结}

AT 牺牲平行度换取稳定的条件生成；NAT 追求平行化，但需要额外处理长度和多模态一致性。实际系统中，NAT 往往要配合蒸馏、迭代修正或其他技巧才更接近 AT 表现。

% =====================================================================
\section{Cross Attention：Encoder 与 Decoder 的桥}
% =====================================================================

Decoder 不只看自己的历史输出，还必须读取 encoder 对输入序列的表示。Transformer decoder 中负责这件事的模块叫 cross attention。原始架构图里，decoder block 有三层子结构：masked multi-head self-attention、cross attention、feed-forward network。Cross attention 位于 masked self-attention 之后。

\begin{figure}[H]
    \centering
    \includegraphics[width=0.78\textwidth]{figures/fig05_cross_attention_position.jpg}
    \caption{原始 Transformer 图中，cross attention 位于 decoder block 中间，是 encoder 和 decoder 之间传递信息的桥梁。\protect\footnotemark}
    \label{fig:cross-position}
\end{figure}
\footnotetext{视频画面时间区间：约 24:00--24:35；字幕对应“cross attention 是连接 encoder 和 decoder 之间的桥梁”“两个输入来自 encoder，一个箭头来自 decoder”。}

\subsection{Q 来自 Decoder，K/V 来自 Encoder}

Cross attention 和 self-attention 的计算形式很像，但 $Q,K,V$ 的来源不同。在 cross attention 中：
\begin{itemize}
    \item query $q$ 来自 decoder 当前状态；
    \item keys $k_i$ 和 values $v_i$ 来自 encoder 输出；
    \item decoder 用 query 去 encoder 输出里“查询”当前需要的信息。
\end{itemize}

具体地，若 encoder 输出为 $h_1,\ldots,h_T$，decoder 当前状态为 $s_t$，则
\[
    q_t = W_Q s_t,\qquad
    k_i = W_K h_i,\qquad
    v_i = W_V h_i.
\]
注意力权重为
\[
    \alpha_{t,i}
    =
    \frac{\exp(q_t^\top k_i/\sqrt{d_k})}
    {\sum_{j=1}^{T}\exp(q_t^\top k_j/\sqrt{d_k})},
\]
输出为
\[
    c_t=\sum_{i=1}^{T}\alpha_{t,i}v_i.
\]
这个 $c_t$ 就是 decoder 在第 $t$ 步从输入序列中抽取出来的上下文信息。

\begin{figure}[H]
    \centering
    \includegraphics[width=0.84\textwidth]{figures/fig06_cross_attention_qkv.jpg}
    \caption{Cross attention 的 $q$ 来自 decoder，$k,v$ 来自 encoder；decoder 通过 query 从 encoder 表示中抽取当前步需要的信息。\protect\footnotemark}
    \label{fig:cross-qkv}
\end{figure}
\footnotetext{视频画面时间区间：约 25:35--26:20；字幕对应“Q 来自 decoder，K 跟 V 来自 encoder”“decoder 凭借产生一个 Q 去 encoder 这边抽取资讯”。}

\subsection{Cross Attention 早于 Transformer}

老师特别指出，cross attention 并不是 Transformer 才有。早期基于 LSTM 的 Seq2seq 语音辨识模型已经使用类似 attention 机制，让 decoder 生成每个输出字符时对齐输入声音片段。Transformer 的新意在于大量使用 self-attention；encoder-decoder 之间的 attention 连接则有更早历史。

\begin{figure}[H]
    \centering
    \includegraphics[width=0.86\textwidth]{figures/fig07_attention_alignment_history.jpg}
    \caption{早期 Listen, Attend and Spell 中已经有输入声音和输出字符之间的 attention 对齐；这不是 Transformer，但说明 cross attention 的思想更早存在。\protect\footnotemark}
    \label{fig:las}
\end{figure}
\footnotetext{视频画面时间区间：约 27:35--28:20；字幕对应“sequence to sequence 用在语音辨识上似乎有潜力”“在有 Transformer 之前已经有 cross attention”。}

\subsection{本章小结}

Cross attention 让 decoder 每一步都能读取 encoder 输出。Self-attention 是同一序列内部互看；cross attention 是 decoder 用 query 去输入序列表示中取信息。Transformer decoder 正是靠 masked self-attention、cross attention 和 FFN 共同完成生成。

% =====================================================================
\section{训练 Decoder：Cross Entropy 与 Teacher Forcing}
% =====================================================================

训练 Seq2seq 时，资料给出输入序列和正确输出序列。例如声音输入对应标签“机器学习”。训练目标是让 decoder 在每个位置输出的分布接近正确 token 的 one-hot vector。对第 $t$ 个位置，若正确 token 是 $y_t^\ast$，损失为
\[
    L_t=-\log P_\theta(y_t^\ast\mid y_{<t}^\ast,H).
\]
整句损失为
\[
    L=\sum_{t=1}^{N+1}L_t,
\]
其中 $N+1$ 通常包含 END token。

\begin{figure}[H]
    \centering
    \includegraphics[width=0.84\textwidth]{figures/fig08_cross_entropy_training.jpg}
    \caption{训练时每个输出分布都和对应的 one-hot ground truth 计算 cross entropy；这和多分类问题本质相同。\protect\footnotemark}
    \label{fig:ce}
\end{figure}
\footnotetext{视频画面时间区间：约 33:25--34:00；字幕对应“正确答案是 one-hot vector”“希望 distribution 跟 one-hot 越接近越好”“计算 cross entropy”。}

\subsection{Teacher Forcing}

训练时常用 teacher forcing：decoder 的输入前缀不是模型自己上一轮猜出的 token，而是正确答案前缀。也就是说，训练第 $t$ 个输出时，模型看到的是
\[
    (\mathrm{BEGIN},y_1^\ast,\ldots,y_{t-1}^\ast),
\]
而不是模型预测的
\[
    (\mathrm{BEGIN},\hat y_1,\ldots,\hat y_{t-1}).
\]

\begin{figure}[H]
    \centering
    \includegraphics[width=0.86\textwidth]{figures/fig09_teacher_forcing.jpg}
    \caption{Teacher forcing 使用正确答案前缀作为 decoder 输入，因此训练可以并行计算每个位置的 cross entropy。\protect\footnotemark}
    \label{fig:teacher-forcing}
\end{figure}
\footnotetext{视频画面时间区间：约 36:55--37:30；字幕对应“把正确的答案当做 decoder 的输入”“训练的时候 decoder 有偷看到正确答案，测试的时候没有”。}

Teacher forcing 的好处是训练稳定且可平行化。因为所有正确前缀都已知，模型可以一次处理整段目标序列，而不用真的一步一步等待自己的输出。但它也埋下一个问题：训练时模型总在干净前缀上学习，测试时却必须面对自己可能产生的错误前缀。

\subsection{本章小结}

Decoder 训练本质上是在每个时间步做分类，用 cross entropy 监督目标 token。Teacher forcing 让训练更容易、更快，但也造成训练和推论输入分布不一致，这会在后面变成 exposure bias。

% =====================================================================
\section{Copy Mechanism 与 Guided Attention}
% =====================================================================

课程接着介绍一些训练 Seq2seq/Transformer 的常见技巧。它们不是 Transformer 架构本身的必备组件，而是针对特定任务的额外归纳偏置。

\subsection{Copy Mechanism}

在对话系统里，如果用户说“你好，我是库洛洛”，机器回复“库洛洛你好，很高兴认识你”，最自然的行为是把用户输入中的名字复制到输出中。若用户说“小杰不能使用念能力了”，机器也可能直接复述“不能使用念能力”这段词组，再追问是什么意思。

\begin{figure}[H]
    \centering
    \includegraphics[width=0.86\textwidth]{figures/fig10_copy_mechanism_chatbot.jpg}
    \caption{Copy mechanism 让模型能把输入中的专有名词或关键短语直接复制到输出，对 chatbot 等任务很有用。\protect\footnotemark}
    \label{fig:copy-chat}
\end{figure}
\footnotetext{视频画面时间区间：约 39:05--39:55；字幕对应“直接把某某某复制出来”“从使用者输入 copy 一些词汇当作输出”。}

摘要任务也有同样需求。很多摘要并不是完全创造新句子，而是从原文复制重要词句，再做压缩和改写。Pointer-generator network 的直觉就是让模型在“从 vocabulary 生成”和“从 source 复制”之间切换：
\[
    P(w)
    =
    p_{\mathrm{gen}}P_{\mathrm{vocab}}(w)
    +
    (1-p_{\mathrm{gen}})
    \sum_{i:x_i=w}\alpha_i.
\]
其中：
\begin{itemize}
    \item $P_{\mathrm{vocab}}(w)$ 是普通生成分布；
    \item $\alpha_i$ 是 source position $i$ 的 attention weight；
    \item $p_{\mathrm{gen}}$ 决定这一刻更像“生成”还是“复制”。
\end{itemize}

\begin{figure}[H]
    \centering
    \includegraphics[width=0.86\textwidth]{figures/fig11_copy_mechanism_summarization.jpg}
    \caption{摘要任务中，许多词汇直接来自原文；pointer-generator 类模型把 vocabulary distribution 和 attention-based copy distribution 混合。\protect\footnotemark}
    \label{fig:copy-summary}
\end{figure}
\footnotetext{视频画面时间区间：约 40:55--41:30；字幕对应“做摘要的时候很多词汇直接从原来的文章里面复制出来”“pointer network”。}

\subsection{Guided Attention}

某些任务中，输入和输出天然单调对齐。例如语音辨识和 TTS 通常从左到右读输入、从左到右产生输出。若 attention 在生成过程中反复跳回去或跳太远，合成语音可能漏字、重复、卡住。Guided attention 就是把这种任务先验写进训练，鼓励 attention 分数沿着单调方向移动。

\begin{figure}[H]
    \centering
    \includegraphics[width=0.88\textwidth]{figures/fig12_guided_attention.jpg}
    \caption{Guided attention 对 attention 形状施加约束：在语音辨识、TTS 等单调对齐任务里，attention 应大致从左到右移动。\protect\footnotemark}
    \label{fig:guided}
\end{figure}
\footnotetext{视频画面时间区间：约 46:35--47:20；字幕对应“attention 有问题会合不出好的结果”“要求机器学到 attention 应该由左向右”。}

\subsection{本章小结}

Copy mechanism 解决“输出应直接复用输入片段”的问题；guided attention 解决“输入输出有已知对齐结构”的问题。它们都体现同一个原则：当任务有明确结构先验时，不必完全交给模型自己摸索。

% =====================================================================
\section{解码策略：Greedy、Beam Search 与 Sampling}
% =====================================================================

训练好 decoder 后，推论时还要决定如何从每一步的分布中选 token。最简单的是 greedy decoding：每一步都选概率最大的 token。但局部最优不等于全局最优。第一步选概率最大的 token，可能导致后续路径整体概率较低。

Beam search 保留多个候选前缀，而不是只保留一个。若 beam size 为 $B$，每一步扩展所有候选，再保留分数最高的 $B$ 个。常见序列分数是 log probability 之和：
\[
    \mathrm{score}(y_{1:t})
    =
    \sum_{i=1}^{t}\log P(y_i\mid y_{<i},H).
\]

\begin{figure}[H]
    \centering
    \includegraphics[width=0.88\textwidth]{figures/fig13_beam_search.jpg}
    \caption{Greedy decoding 只走红色局部最优路径；beam search 用近似搜索保留多条候选，尝试接近绿色全局较优路径。\protect\footnotemark}
    \label{fig:beam}
\end{figure}
\footnotetext{视频画面时间区间：约 50:35--51:15；字幕对应“无法穷举所有路径”“每个地方分叉都是很多可能”“beam search 找 approximate solution”。}

\subsection{什么时候 Beam Search 有用}

Beam search 对答案明确的任务更有帮助，例如语音辨识、机器翻译等。目标通常只有少数合理答案，找到高概率序列往往能提升结果。对开放式生成、故事续写、聊天等任务，过强的 beam search 可能产生无聊、重复、模板化输出，因为“最可能的句子”未必是最有趣或最自然的句子。

\subsection{Sampling}

当任务需要多样性时，可以从分布中随机采样，而不是总取最大值。Sampling 可以让生成结果更丰富，但也会引入不稳定。实际系统常配合 temperature、top-k、top-p 等方法控制随机性。

\begin{figure}[H]
    \centering
    \includegraphics[width=0.88\textwidth]{figures/fig14_sampling_generation.jpg}
    \caption{开放式文本生成中，beam search 可能过于保守；适度随机性有助于产生更自然、多样的输出。\protect\footnotemark}
    \label{fig:sampling}
\end{figure}
\footnotetext{视频画面时间区间：约 53:10--53:40；字幕对应“需要机器发挥一点创造力”“不是只有一个答案的任务”。}

\subsection{评价指标与强化学习}

老师还提到：如果你真正关心的是 BLEU 等句子级评价指标，逐 token 的 cross entropy 未必直接优化这些指标。这时可以考虑 reinforcement learning，把不可微或句子级指标当作 reward 来优化。不过这条路通常更难训练，也需要谨慎设计 reward，否则模型可能学会钻指标漏洞。

\begin{warningbox}{训练目标和最终指标可能不一致}
    Cross entropy 优化的是下一个 token 的条件概率；BLEU、ROUGE、人类偏好等指标评价的是整段输出。两者相关但不相同。许多生成模型问题都来自这个错位：局部概率高，不代表整体结果好。
\end{warningbox}

\subsection{本章小结}

Greedy 简单但短视；beam search 是近似全局搜索，适合答案较明确的任务；sampling 增加多样性，适合开放式生成。解码策略不是训练后的细枝末节，而会显著影响最终表现。

% =====================================================================
\section{Exposure Bias 与 Scheduled Sampling}
% =====================================================================

Teacher forcing 带来训练和推论不一致。训练时 decoder 总是看到正确前缀；推论时 decoder 看到自己的预测前缀。若模型从没在训练时见过错误前缀，测试时一旦早期出错，就可能一步错、步步错。这个现象叫 exposure bias。

\begin{figure}[H]
    \centering
    \includegraphics[width=0.86\textwidth]{figures/fig15_exposure_bias.jpg}
    \caption{Exposure bias：训练时 decoder 输入是正确前缀，测试时输入却来自模型自己的预测；一旦预测错，后续状态就偏离训练分布。\protect\footnotemark}
    \label{fig:exposure}
\end{figure}
\footnotetext{视频画面时间区间：约 58:50--59:20；字幕对应“训练的时候都是正确的，测试的时候模型看到自己的输出”“这个不一致叫 exposure bias”。}

\subsection{Scheduled Sampling 的直觉}

一种直觉解法是：训练时不要永远喂正确 token，偶尔把模型自己预测的 token 喂回去。这样模型会在训练中见到一些错误前缀，学会从不完美状态恢复。这个方法叫 scheduled sampling。

\begin{figure}[H]
    \centering
    \includegraphics[width=0.86\textwidth]{figures/fig16_scheduled_sampling.jpg}
    \caption{Scheduled sampling 让训练时的 decoder 偶尔接收模型采样输出，试图缓解 teacher forcing 和推论之间的 mismatch。\protect\footnotemark}
    \label{fig:scheduled}
\end{figure}
\footnotetext{视频画面时间区间：约 59:35--60:15；字幕对应“偶尔给它一些错的东西”“这一招叫 scheduled sampling”“对 Transformer 来说另有招数”。}

不过传统 scheduled sampling 会破坏 Transformer 的平行训练能力。因为一旦某个位置输入依赖模型上一位置的实际采样结果，训练就不能简单地一次并行算完所有位置。因此，Transformer 里的 scheduled sampling 需要专门设计，不能直接照搬早期 RNN/LSTM 的做法。

\begin{warningbox}{Scheduled Sampling 不是学习率排程}
    课程中特别提醒：scheduled sampling 和 learning rate schedule 是两件完全不同的事。前者改变 decoder 训练时看到的输入 token；后者改变优化器更新参数的步长。
\end{warningbox}

\subsection{本章小结}

Exposure bias 来自 teacher forcing：训练看正确前缀，推论看模型前缀。Scheduled sampling 的思想是让训练更像推论，但在 Transformer 中会牵涉平行化和训练稳定性，需要额外设计。

% =====================================================================
\section{总结与延伸}
% =====================================================================

本讲把 Transformer 的 decoder 和训练技巧补完整。若把两讲合起来看，Transformer 的完整结构可以概括为：
\[
\begin{aligned}
    \text{Encoder}
    &: \mathrm{Embed+Pos}
    \rightarrow
    (\mathrm{SelfAttn+FFN})^{N},\\
    \text{Decoder}
    &: \mathrm{ShiftedEmbed+Pos}
    \rightarrow
    (\mathrm{MaskedSA+CrossAttn+FFN})^{N}.
\end{aligned}
\]

Decoder 这半边最重要的概念有五个：
\begin{itemize}
    \item Autoregressive generation：一步一步产生 token；
    \item Masked self-attention：不能偷看未来 token；
    \item Cross attention：用 decoder query 读取 encoder key/value；
    \item Teacher forcing：训练时用正确前缀，换取稳定和并行；
    \item Exposure bias：训练/推论前缀分布不一致，可能造成错误传播。
\end{itemize}

\begin{importantbox}{本讲最重要的理解}
    Transformer decoder 是一个条件生成器。它既要遵守因果顺序，又要读取输入序列，还要在训练时被逐 token 监督。很多“技巧”不是装饰，而是在补 decoder 的真实缺口：复制输入、保持单调对齐、近似搜索、增加随机性、缓解 exposure bias。
\end{importantbox}

从工程角度看，Transformer 不是一张架构图就结束了。训练目标、teacher forcing、解码策略、停止条件和任务先验都会强烈影响最终系统表现。学会 encoder/decoder 的方块连接只是第一步；真正能用好 Seq2seq 模型，还要理解生成过程中的这些细节。

\subsection{拓展阅读}

\begin{itemize}
    \item \href{https://arxiv.org/abs/1706.03762}{Attention Is All You Need}：原始 Transformer；
    \item \href{https://arxiv.org/abs/1609.08144}{Listen, Attend and Spell}：早期语音辨识中的 encoder-decoder attention；
    \item \href{https://arxiv.org/abs/1704.04368}{Pointer-generator Network}：摘要中的 copy mechanism；
    \item \href{https://arxiv.org/abs/1904.09751}{The Curious Case of Neural Text Degeneration}：beam search 与 sampling 的生成质量问题；
    \item \href{https://arxiv.org/abs/1506.03099}{Scheduled Sampling}：缓解 exposure bias 的早期方法。
\end{itemize}

\end{document}
